\chapter{От статического селектора к динамике переноса дефекта}

Первая часть Тома V последовательно проверяла, может ли одна родительская
архитектура выбрать наблюдаемый сектор посредством положительного следа,
кривизны, отображения момента, связности или спектрального кручения. Эти
проверки дали содержательную кинематику, но не вывели единственный
ненулевой оператор Юкавы. Настоящая глава не отменяет полученные запреты.
Она меняет вопрос: вместо выбора готового состояния исследуется закон
локального переноса, благодаря которому устойчивое возбуждение или дефект
может существовать как распространяющийся процесс.

\section{Обратная сверка с исходной программой}

В ранних текстах проекта вакуум описывался как упорядоченная связная среда,
частица --- как устойчивый узел её организации, а нейтринный сектор --- как
подвижный геометрический дефект. Эти формулировки содержатся, в частности,
в \path{s2t/RPFT-main/habr/print2.md} и в архивной проработке
\path{архив-2025-2026/2026-02-проработка/Проработка/interpritate.md}.
Тогда это была онтологическая гипотеза без достаточного операторного языка.

Позднейшие тома независимо построили четыре части такого языка:
\begin{enumerate}
 \item ранг-один проектор как устойчивую опорную конфигурацию;
 \item голономию как сохраняемое чтение обхода вокруг дефекта;
 \item соответствие Мориты как пространство допустимых переходов между
 двумя операторными чтениями;
 \item связность и кручение как средства сравнения соседних конфигураций.
\end{enumerate}
Тем самым ранняя картина не становится доказанной физикой, но получает
строгую переформулировку: первичным кандидатом является не вещество внутри
заранее заданного пространства, а локальный закон согласованной передачи
операторной конфигурации.

\section{Граница с известной математической физикой}

Переход от образа к проверяемой гипотезе требует отделить известные факты
от проектной интерпретации.

В геометрической теории дефектов дислокации описываются кручением
римановой--картановой геометрии, а дисклинации --- кривизной; см.
\path{arXiv:cond-mat/0407469} и \path{arXiv:0909.4068}. Локальная,
однородная и унитарная динамика на дискретном носителе может иметь в
непрерывном пределе уравнение Дирака; см. \path{arXiv:1307.3524} и
\path{arXiv:1212.2839}. Топологический индекс способен защищать
майорановскую нулевую моду, связанную с вихрем или линейным дефектом; см.
\path{arXiv:0911.2558} и \path{arXiv:1003.4814}. Квантовые блуждания также
могут воспроизводить стандартную кинематику нейтринных осцилляций; см.
\path{arXiv:1607.00529}.

Ни один из этих результатов не доказывает, что физическое нейтрино является
дислокацией вакуума, а наблюдаемые фермионы --- узлами одной фундаментальной
нити. Кроме того, спектральное кручение конечной геометрии, вычисленное в
предыдущей главе, нельзя без дополнительного словаря отождествлять с
кручением римановой--картановой среды. Настоящая ветвь принимает только
более слабую гипотезу: проектные операторы могут реализовать ту же
математическую схему \emph{локальный дефект --- перенос --- непрерывный
дираковский предел}.

\section{Что именно переосмысляется}

\begin{definition}[Статический селектор]
Статическим селектором называется функционал $V(q)$ на пространстве
конфигураций, физическое содержание которого связывается с выделением
единственного минимума или единственной орбиты $q_*$.
\end{definition}

\begin{definition}[Динамический перенос]
Динамическим переносом называется локальное отображение
\begin{equation}
 U_e:\mathcal H_x\longrightarrow\mathcal H_y,
 \qquad e=(x,y),
 \label{eq:v5-local-transfer}
\end{equation}
согласованное с алгебрами концов, вещественной структурой, градуировкой и
общим скалярным произведением. Физическое возбуждение задаётся не одной
точкой конфигурационного пространства, а устойчивым классом орбит такого
переноса.
\end{definition}

Провал статического селектора не запрещает динамический перенос. Например,
равномерная среда может иметь непрерывное семейство эквивалентных положений
дефекта; наличие нулевой моды тогда выражает возможность поступательного
движения, а не физическую недоопределённость внутреннего строения. Поэтому
нулевые направления гессиана следует разделять на три класса:
\begin{enumerate}
 \item калибровочные направления, удаляемые факторизацией;
 \item моды положения или фазы устойчивого дефекта;
 \item настоящие невыведенные внутренние параметры.
\end{enumerate}
Предыдущие гейты часто фиксировали число нулевых направлений, но не всегда
могли выполнить это динамическое различение, поскольку оператор переноса
не был построен.

\section{Новый словарь объектов проекта}

Следующая таблица является не доказанным отождествлением, а рабочим
словарём, который должен быть проверен одним операторным построением:
\begin{center}
\begin{tabular}{p{0.30\textwidth}p{0.61\textwidth}}
\toprule
Раннее понятие & Операторный кандидат \\
\midrule
упорядоченная среда & сеть локальных алгебр, модулей и согласованных
соответствий \\
устойчивый узел & конечная по энергии орбита проекторов или связностей с
нетривиальным индексом \\
перемещение узла & унитарная или изометрическая передача локального класса
между соседними углами \\
дислокация & дефект композиции переносов, измеряемый кручением или
трансляционной голономией \\
дисклинация & дефект вращательной голономии, измеряемый кривизной \\
заряд & дискретный индекс или класс голономии замкнутого обхода \\
масса & щель дисперсионного оператора либо цена смешивания встречных
хиральных каналов \\
аромат & внутренний канал чтения одного класса переноса \\
осцилляция & унитарная эволюция между каналами чтения при сохранении
полного дефектного класса \\
\bottomrule
\end{tabular}
\end{center}

Особенно существенно новое чтение носителя
\begin{equation}
 E=M_{20\times15}(\mathbb C).
\end{equation}
Он не обязан перечислять $300$ частиц или независимых физических полей.
Он может быть пространством элементарных способов перехода между
семейным и наблюдаемым чтениями. Это согласуется с уже установленным
статусом $E$ как соответствия Мориты и объясняет, почему попытка
калибровать всю связывающую алгебру $M_{35}(\mathbb C)$ была ошибочной:
пространство допустимых переходов не тождественно координатной алгебре
наблюдаемого мира.

\section{Минимальная математическая постановка переноса}

Пусть $\Gamma$ --- пока не физическая решётка пространства, а конечный
граф локальных сравнений. В каждой вершине $x$ заданы копии уже выведенных
данных
\begin{equation}
 (\mathcal A_x,\mathcal H_x,J_x,\gamma_x,\rho_x,\tau_x).
\end{equation}
Ребру $e=(x,y)$ сопоставляется ограниченное соответствие $E_e$ и связность
$\nabla_e$. Элементарный перенос должен иметь вид
\begin{equation}
 U_e=\mathcal P\exp\!\left(-\int_e\nabla\right)
 :\rho_x\mathcal H_x\longrightarrow\rho_y\mathcal H_y
 \label{eq:v5-connection-transfer}
\end{equation}
или его конечномерный аналог через частичную изометрию полярного
разложения. Композиция вдоль пути $p=e_n\cdots e_1$ равна
\begin{equation}
 U(p)=U_{e_n}\cdots U_{e_1}.
\end{equation}
Замкнутый путь определяет голономию, а нарушение независимости от
стягиваемой деформации пути измеряет кривизну. Дефект трансляционного типа
должен дополнительно проявляться в некоммутативности дискретного переноса с
оператором локальной степени или высоты.

Глобальный одношаговый оператор записывается как
\begin{equation}
 W=\sum_{e\in\Gamma_1}S_e\otimes U_e,
 \label{eq:v5-global-step}
\end{equation}
где $S_e$ переносит опору между соседними вершинами. Он допустим только при
выполнении следующих условий:
\begin{align}
 W^*W&=WW^*=I,
 \label{eq:v5-transfer-unitarity}\\
 [W,a]&\ \text{локален},\\
 JW&=WJ,\qquad \gamma W\gamma=W^*
 \quad\text{либо выполняется эквивалентное условие},\\
 \tau(W^*XW)&=\tau(X).
\end{align}
Здесь первая строка выражает унитарность, вторая --- конечную скорость
передачи, третья --- вещественность и хиральность, а четвёртая ---
сохранение единственной меры.
Точная форма условия с $\gamma$ должна следовать из выбранной дискретизации,
а не подбираться после вычисления спектра.

Для однородной ветви выполняется преобразование Фурье и определяется
эффективный генератор
\begin{equation}
 W(k)=\exp[-iaH_{\rm eff}(k)].
\end{equation}
Критическим тестом является разложение
\begin{equation}
 H_{\rm eff}(k)=c\,\alpha^ik_i+\beta M+O(a|k|^2),
 \label{eq:v5-dirac-limit}
\end{equation}
где $c$ должен быть общим для всех наблюдаемых каналов, а матрица $M$
должна возникать из уже имеющейся геометрии переноса. Если в
\eqref{eq:v5-dirac-limit} приходится независимо задавать скорости,
переходные амплитуды или массы для $u,d,e$, то статическая неоднозначность
только переименована и прогресса нет.

\section{Масса как щель переноса}

В статическом подходе масса искалась как кривизна потенциала в выбранном
минимуме. Динамическая постановка допускает более первичное определение.

\begin{definition}[Переносная масса]
Для ветви дисперсии $\omega_s(k)$ переносной массой называется щель
\begin{equation}
 m_s=\frac1{c^2}\min_k|\omega_s(k)|,
 \end{equation}
если низкоэнергетическое разложение имеет релятивистский вид
\begin{equation}
 \omega_s(k)^2=c^2|k|^2+m_s^2c^4+O(a|k|^3).
 \label{eq:v5-relativistic-dispersion}
\end{equation}
\end{definition}

Такое определение превращает вопрос о юкавских отношениях в вопрос о
спектре одного оператора переноса. Однако оно ничего не решает само по
себе: произвольные межканальные амплитуды дадут произвольные массы. Новый
результат появится только в том случае, если относительные щели будут
жёстко следовать из ранга проектора, тетраэдрической голономии,
двусторонней связности и единственного следа.

\section{Нейтрино как отдельный класс переноса}

Развилка $H_{15}/H_{16}$ теперь получает содержательное толкование.
Заряженные поля $u,d,e$ принадлежат одноформенному сектору $H_{15}$, тогда
как нейтринная масса не обязана быть четвёртым ребром того же типа. Она
может принадлежать комплексу высшей степени
\begin{equation}
 \mathcal C^0\xrightarrow{d_0}\mathcal C^1
 \xrightarrow{d_1}\mathcal C^2,
 \qquad d_1d_0=0,
 \label{eq:v5-neutrino-complex}
\end{equation}
где физический перенос действует на классе когомологий или на
локализованной нулевой моде дефекта.

Это согласуется с ранее полученным условным майорановским результатом:
нетривиальный вихревой класс способен защищать одну вещественную нулевую
моду после возникновения спаривающего конденсата. Но отсутствующий вывод
самого конденсата сохраняется. Поэтому выражение «нейтрино есть бегущая
дислокация» пока является названием исследовательской гипотезы, а не
теоремой.

Три аромата допустимо интерпретировать как три проектора чтения
$P_e,P_\mu,P_\tau$ одного переносимого внутреннего состояния. Тогда
вероятность перехода имела бы стандартную операторную форму
\begin{equation}
 \mathcal P_{\alpha\to\beta}(n)
 =\left\|P_\beta W^nP_\alpha\psi\right\|^2.
 \label{eq:v5-flavour-transfer}
\end{equation}
Но проект должен вывести сами проекторы, спектральные фазы и их связь с
семейной геометрией. Простое воспроизведение известной матрицы смешивания
не будет считаться объяснением.

\section{Что сохраняется из первой части}

Новая постановка не разрешает возвращаться к закрытым приёмам. Сохраняются
следующие результаты:
\begin{enumerate}
 \item обычный спектр $D^2$ не различает ориентированную разность матриц
 Грама;
 \item высота, скручивание или центральный уровень не могут вводиться как
 скрытый внешний селектор;
 \item положительная статическая норма связности выбирает нуль;
 \item спектральное кручение измеряет уже заданные юкавские отношения, но
 не выбирает их;
 \item $M_{35}(\mathbb C)$ является связывающим контейнером, а не новой
 физической калибровочной алгеброй;
 \item условная майорановская мода не доказывает динамического рождения
 дефекта.
\end{enumerate}
Следовательно, часть II не является обходом отрицательных результатов. Она
использует их для запрета произвольных коэффициентов в операторе переноса.

\section{Следующий решающий гейт}

Следующий гейт должен проверить \emph{локальный оператор переноса
дефектного класса}. Его минимальная последовательность такова:
\begin{enumerate}
 \item выбрать простейший граф сравнений, на котором уже определены
 $\rho$, $E=M_{20\times15}(\mathbb C)$ и двусторонняя связность;
 \item классифицировать все локальные частичные изометрии $U_e$,
 совместимые с двумя действиями алгебр, KO6, градуировкой и следом;
 \item удалить только унитарно эквивалентные базисные выборы, не устраняя
 физические нулевые моды объявлением новой калибровочной группы;
 \item проверить, остаётся ли единственная орбита одношагового оператора
 $W$ с точностью до общего масштаба шага;
 \item вычислить точный спектр $W(k)$ и проверить дираковский предел
 \eqref{eq:v5-dirac-limit}, общий предельный световой конус и отсутствие
 фермионного удвоения либо его физическое проектирование;
 \item выделить топологический дефект и проверить, переносится ли его индекс
 без переноса локального вещества носителя;
 \item только после этого вычислять щели, хиральность и возможные
 вероятности ароматных переходов.
\end{enumerate}

\begin{nogocriterion}[Запрет переименованной подгонки]
Ветка закрывается, если наиболее общий допустимый $W$ содержит независимые
амплитуды переноса для физических секторов, числом не меньшие уже известных
юкавских свобод, либо если требуемые массы и смешивания возникают только
после их подстановки в локальные матрицы $U_e$.
\end{nogocriterion}

\begin{nogocriterion}[Запрет буквальной кристаллической метафоры]
Совпадение языка дефектов, кручения и голономии недостаточно. Ветка
закрывается как физическая интерпретация, если она не выводит унитарность,
локальность, общий релятивистский предел, наблюдаемую хиральность и
совместимость с калибровочной группой Стандартной модели. Механическая
решётка с выделенной системой покоя не считается допустимым результатом.
\end{nogocriterion}

\begin{theorem}[Статус прозрения]
Совокупность результатов проекта допускает непротиворечивое новое чтение:
родитель Мориты является кинематикой переходов, связность задаёт правило их
сравнения, голономия и индекс классифицируют устойчивый дефект, а масса
может быть спектральной щелью его переноса. Это чтение не выводит
наблюдаемую динамику, но формулирует ранее отсутствовавший решающий объект
--- локальный унитарный оператор переноса --- и запрещает возвращение к
произвольному статическому потенциалу.
\end{theorem}

\begin{protocol}
\begin{enumerate}
 \item связь с ранней гипотезой нити и дефекта:
 \textbf{восстановлена как проверяемый операторный словарь};
 \item буквальная модель механического эфира:
 \textbf{не принимается};
 \item родитель Мориты как пространство переходов:
 \textbf{сохранён};
 \item статический вывод юкавских коэффициентов:
 \textbf{по-прежнему не получен};
 \item масса как возможная щель переноса:
 \textbf{сформулирована, но не вычислена};
 \item нейтрино как дефектный перенос высшей степени:
 \textbf{допустимая гипотеза, не теорема};
 \item следующий объект:
 \textbf{локальный унитарный оператор переноса};
 \item следующий критерий успеха:
 \textbf{единственная орбита, дираковский предел и спектр без подгонки};
 \item физическое замыкание:
 \textbf{не достигнуто}.
\end{enumerate}
\end{protocol}